\documentclass[UTF8,11pt]{ctexart}
\usepackage[a4paper,margin=24mm]{geometry}
\usepackage{amsmath,amssymb,amsthm,mathtools,mathrsfs}
\usepackage[all]{xy}
\usepackage{xurl,hyperref,enumitem}
\usepackage{tikz}
\usetikzlibrary{arrows.meta,cd,decorations.markings,calc,patterns}
\hypersetup{hidelinks,breaklinks=true}
\setlength{\emergencystretch}{3em}
\allowdisplaybreaks
\newtheorem*{theorem}{Theorem}
\newtheorem*{thm}{Theorem}
\newtheorem*{lemma}{Lemma}
\newtheorem*{proposition}{Proposition}
\newtheorem*{prop}{Proposition}
\newtheorem*{corollary}{Corollary}
\newtheorem*{cor}{Corollary}
\newtheorem*{claim}{Claim}
\theoremstyle{definition}
\newtheorem*{definition}{Definition}
\newtheorem*{defn}{Definition}
\newtheorem*{remark}{Remark}
\newtheorem*{example}{Example}
\newcommand{\R}{\mathbb R}
\newcommand{\C}{\mathbb C}
\newcommand{\Z}{\mathbb Z}
\newcommand{\Q}{\mathbb Q}
\newcommand{\N}{\mathbb N}
\newcommand{\F}{\mathbb F}
\newcommand{\D}{\mathbb D}
\newcommand{\bB}{\mathbb B}
\newcommand{\bC}{\mathbb C}
\newcommand{\bR}{\mathbb R}
\newcommand{\bZ}{\mathbb Z}
\newcommand{\bN}{\mathbb N}
\newcommand{\bQ}{\mathbb Q}
\newcommand{\bD}{\mathbb D}
\newcommand{\bF}{\mathbb F}
\newcommand{\CP}{\mathbb{CP}}
\newcommand{\RP}{\mathbb{RP}}
\newcommand{\sO}{\mathscr O}
\newcommand{\sI}{\mathscr I}
\newcommand{\sF}{\mathscr F}
\newcommand{\sG}{\mathscr G}
\newcommand{\sH}{\mathscr H}
\newcommand{\sR}{\mathscr R}
\newcommand{\sM}{\mathscr M}
\newcommand{\sV}{\mathscr V}
\newcommand{\sU}{\mathscr U}
\DeclareMathOperator{\Hom}{Hom}
\DeclareMathOperator{\Ext}{Ext}
\DeclareMathOperator{\Tor}{Tor}
\DeclareMathOperator{\Spec}{Spec}
\DeclareMathOperator{\Proj}{Proj}
\DeclareMathOperator{\supp}{supp}
\DeclareMathOperator{\im}{im}
\DeclareMathOperator{\id}{id}
\DeclareMathOperator{\Tr}{Tr}
\DeclareMathOperator{\tr}{tr}
\DeclareMathOperator{\rank}{rank}
\DeclareMathOperator{\ord}{ord}
\DeclareMathOperator{\dist}{dist}
\DeclareMathOperator{\End}{End}
\DeclareMathOperator{\Aut}{Aut}
\DeclareMathOperator{\Gal}{Gal}
\DeclareMathOperator{\vol}{vol}
\DeclareMathOperator{\Cl}{Cl}
\DeclareMathOperator{\coker}{coker}
\DeclareMathOperator{\codim}{codim}
\DeclareMathOperator{\divergence}{div}
\DeclareMathOperator{\Ric}{Ric}
\DeclareMathOperator{\grad}{grad}
\DeclareMathOperator{\Hess}{Hess}
\newcommand{\abs}[1]{\left\lvert#1\right\rvert}
\newcommand{\sabs}[1]{\lvert#1\rvert}
\newcommand{\babs}[1]{\bigl\lvert#1\bigr\rvert}
\newcommand{\Babs}[1]{\Bigl\lvert#1\Bigr\rvert}
\newcommand{\bbabs}[1]{\biggl\lvert#1\biggr\rvert}
\newcommand{\BBabs}[1]{\Biggl\lvert#1\Biggr\rvert}
\newcommand{\norm}[1]{\left\lVert#1\right\rVert}
\newcommand{\snorm}[1]{\lVert#1\rVert}
\newcommand{\bnorm}[1]{\bigl\lVert#1\bigr\rVert}
\newcommand{\Bnorm}[1]{\Bigl\lVert#1\Bigr\rVert}
\newcommand{\bbnorm}[1]{\biggl\lVert#1\biggr\rVert}
\newcommand{\BBnorm}[1]{\Biggl\lVert#1\Biggr\rVert}
\newcommand{\widebar}[1]{\overline{#1}}
\newcommand{\nicefrac}[2]{\frac{#1}{#2}}
\newcommand{\linnprod}[2]{\langle#1,#2\rangle}
\newcommand{\myindex}[1]{#1}
\newcommand{\glsadd}[1]{}
\newcommand{\avoidbreak}{}
\newcommand{\sourceref}[1]{\textnormal{[原文：\nolinkurl{#1}]}}
\newcommand{\thmref}[1]{\sourceref{#1}}
\newcommand{\lemmaref}[1]{\sourceref{#1}}
\newcommand{\propref}[1]{\sourceref{#1}}
\newcommand{\corref}[1]{\sourceref{#1}}
\newcommand{\defnref}[1]{\sourceref{#1}}
\newcommand{\exampleref}[1]{\sourceref{#1}}
\newcommand{\exerciseref}[1]{\sourceref{#1}}
\newcommand{\figureref}[1]{\sourceref{#1}}
\newcommand{\sectionref}[1]{\sourceref{#1}}
\newcommand{\appendixref}[1]{\sourceref{#1}}
\newcommand{\stacksref}[2]{\href{https://stacks.math.columbia.edu/tag/#1}{#2 [#1]}}


\newcommand{\E}{\mathbb E}
\newcommand{\Prob}{\mathbb P}
\newcommand{\Reff}{\mathcal R_{\mathrm{eff}}}
\newcommand{\eps}{\varepsilon}
\DeclareMathOperator{\diam}{diam}
\DeclareMathOperator{\Var}{Var}
\DeclareMathOperator{\Crit}{Crit}
\newcommand{\dd}{\,\mathrm d}

\begin{document}
\section*{054\quad 有限平面图局部极限的常返性}
\subsection*{来源与计数目标}
Asaf Nachmias, \emph{Planar Maps, Random Walks and Circle Packing}, LNM 2243, Springer, 2020，第5章 \emph{Planar Local Graph Limits}，pp.61--71：Definition 5.1、Theorem 5.2、\S5.2--5.3 的全部核心证明；附第4章 Lemma 4.9。获取日期2026-09-28。
\par\url{https://link.springer.com/chapter/10.1007/978-3-030-27968-4_5}
\par\textbf{本文件只计一个问题：}证明有界度有限平面图的局部极限几乎必然常返；Theorem 5.8 的电阻定量结论及支撑引理合在本题，不另计数。
\subsection*{作者命题与人类证明}
\paragraph{Definition 5.1: Local convergence.}
\noindent\textit{以下背景与记号据所列原文章节摘编；不是逐字引文，也不是新增证明。}\par
Let $\{G_n\}$ be a sequence of (possibly random) finite graphs. We say that $G_n$ converges locally to a (possibly infinite) random rooted graph $(U,\rho)\in\mathcal G_\bullet$, and denote it by $G_n\xrightarrow{\mathrm{loc}}(U,\rho)$, if for every integer $r\ge1$,
\[
B_{G_n}(\rho_n,r)\xrightarrow{d}B_U(\rho,r),
\]
where $\rho_n$ is a uniformly chosen vertex from $G_n$. Here $B_G(\rho,r)$ is the subgraph rooted at $\rho$ spanned by vertices of graph distance at most $r$; rooted graphs are viewed up to root-preserving graph isomorphism.

\begin{theorem}[5.2, Benjamini--Schramm]
Let $M<\infty$ and let $G_n$ be finite planar maps (possibly random) with degrees almost surely bounded by $M$ such that $G_n\xrightarrow{\mathrm{loc}}(U,\rho)$. Then $(U,\rho)$ is almost surely recurrent.
\end{theorem}

\paragraph{The Magic Lemma (Section 5.2).}
Suppose $C\subseteq\R^2$ is finite. For each $w\in C$, define $\rho_w=\min\{|v-w|:v\in C\setminus\{w\}\}$. We call $\rho_w$ the isolation radius of $w$. Given $\delta\in(0,1)$, $s\ge2$ and $w\in C$, we say that $w$ is $(\delta,s)$-supported if
\[
\inf_{p\in\R^2}\left|C\cap B(w,\delta^{-1}\rho_w)\setminus B(p,\delta\rho_w)\right|\ge s.
\]
\begin{lemma}[5.3]
There exists $A>0$ such that for every $\delta\in(0,1/2)$, every finite $C\subseteq\R^2$ and every $s\ge2$, the number of $(\delta,s)$-supported points in $C$ is at most $A|C|\delta^{-2}\ln(\delta^{-1})/s$.
\end{lemma}
\paragraph{Proof of Lemma 5.3.}
Let $k\ge3$ be an integer (later we will take $k=k(\delta)$). Let $G_0$ be a tiling of $\R^2$ by $1\times1$ squares, rooted at some point $p$, and for every $n\in\Z$, let $G_n$ be a tiling of $\R^2$ by $k^n\times k^n$ such that each square of $G_n$ is tiled by $k\times k$ squares of $G_{n-1}$. We may choose $p$ so that none of the points of $C$ lies on the edge of a square.
We say that a square $S\in G_n$ is $s$-supported if for every smaller square $S'\in G_{n-1}$ we have that $|C\cap(S\setminus S')|\ge s$.
\begin{claim}[5.5]
For any $s\ge2$ the total number of $s$-supported squares, in $G=\bigcup_{n\in\Z}G_n$, is at most $2|C|/s$.
\end{claim}
\begin{proof}
Define a ``flow'' $f:G\times G\to\R$ as follows:
\[
f(S',S)=\begin{cases}
\min(s/2,|S'\cap C|)&S'\subseteq S,\ S'\in G_n,\ S\in G_{n+1},\\
-f(S,S')&S\subseteq S',\ S\in G_n,\ S'\in G_{n+1},\\
0&\text{otherwise}.
\end{cases}
\]
Let us make two initial observations. First we have that
\[
\sum_{S'\in G}f(S',S)\ge0,\tag{5.1}
\]
by splitting into the two cases depending on whether there exists a square $S'\subseteq S$ such that $f(S',S)=s/2$ or not. Secondly, if $S$ is a $s$-supported square
\[
\sum_{S'\in G}f(S',S)\ge s/2,\tag{5.2}
\]
by splitting into cases depending on whether the number of squares $S'\subseteq S$ such that $f(S',S)=s/2$ is at most one or at least two.
Let $a\in\Z$ be such that each square in $G_a$ contains at most $1$ point of $C$ so there are no $s$-supported squares in $\bigcup_{n\le a}G_n$. It easily follows from the definition of $f$ that
\[
\sum_{S'\in G_a}\sum_{S\in G_{a+1}}f(S',S)=|C|,\tag{5.3}
\]
and that for every $b\in\Z$
\[
\sum_{S'\in G_b}\sum_{S\in G_{b+1}}f(S',S)\ge0.\tag{5.4}
\]
Now, using (5.3) and (5.4),
\begin{align*}
\sum_{n=a+1}^b\sum_{S\in G_n}\sum_{S'\in G}f(S',S)
&=\sum_{n=a+1}^b\sum_{S\in G_n}\left(\sum_{S'\in G_{n-1}}f(S',S)+\sum_{S'\in G_{n+1}}f(S',S)\right)\\
&=\sum_{S\in G_{a+1}}\sum_{S'\in G_a}f(S',S)
+\sum_{S\in G_b}\sum_{S'\in G_{b+1}}f(S',S)\le|C|.
\end{align*}
Therefore, using (5.1) and (5.2), we deduce that there are at most $2|C|/s$ squares in $\bigcup_{b\ge n>a}G_n$ that are $s$-supported. Sending $b\to\infty$ finishes the proof.
\end{proof}
The above claim is very close to the statement of Lemma 5.3 which we are pursuing. However, we need to move from squares to circles. We choose $k=\lceil20\delta^{-2}\rceil$ and let $\beta\sim\operatorname{Unif}([0,\ln k])$. Let $G_0$ be a tiling with side length $e^\beta$, based at the origin.
Suppose we have defined $G_n$ as a tiling of squares of side length $e^\beta k^n$; then $G_{n+1}$ is a tiling of squares of side length $e^\beta k^{n+1}$ that is based uniformly at one of the $k^2$ possible points of $G_n$. Because the desired statement is invariant under translation and dilation of $C$, we may assume that $C$ does not intersect the edges of $G_n$ (for every $n$) and that $\rho_w\ge k$ for every $w\in C$. We call a point $w\in C$ a city in a square $S\in G$ if the side length of $S$ is in the interval $[4\delta^{-1}\rho_w,5\delta^{-1}\rho_w]$, and the distance from $w$ to the center of $S$ is at most $\delta^{-1}\rho_w$.
\begin{claim}[5.6]
The probability that any given $w\in C$ is a city is $\Omega(\ln^{-1}(\delta^{-1}))$.
\end{claim}
\begin{proof}
For the first item to hold, $\beta$ needs to satisfy that there exists $n\in\Z$ such that $e^\beta k^n\in[4\delta^{-1}\rho_w,5\delta^{-1}\rho_w]$, or $\beta+n\ln k\in\ln(\delta^{-1}\rho_w)+[\ln4,\ln5]$. Since $\beta\in\operatorname{Unif}([0,\ln k])$, the probability for that is $(\ln(5/4))/\ln k$, which is $\Omega(\ln^{-1}(\delta^{-1}))$ when $\delta\in(0,1/2)$.
As for the second item, it holds with positive probability (independent of $\delta$) over the $k^2$ choices for basing $G_n$ on top of $G_{n-1}$, given that $\beta$ satisfies the requirement posed by the first item.
\end{proof}
\begin{claim}[5.7]
If $w$ is a city in $S$ and is $(\delta,s)$-supported, then $S$ is $s$-supported.
\end{claim}
\begin{proof}
If $S\in G_n$ is as above, then any little square $S'\in G_{n-1}$ has side length at most
$\frac{\delta^2}{20}\frac{5\rho_w}{\delta}=\frac{\delta\rho_w}{4}$.
Hence, it is contained in a disk of radius $\delta\rho_w$. Thus, for every $S'\in G_{n-1}$ with $S'\subseteq S$ there exists a point $p$ such that
\[
|C\cap(S\setminus S')|\ge\left|C\cap\bigl(B(w,\delta^{-1}\rho_w)\setminus B(p,\delta\rho_w)\bigr)\right|\ge s,
\]
where we have used the fact that $B(w,\delta^{-1}\rho_w)\subset S$.
\end{proof}
Now note that the expected number of pairs $(w,S)$ such that $S$ is $s$-supported, $w$ is $(\delta,s)$-supported, and $w$ is a city, is at least $c\ln^{-1}(\delta^{-1})N$, where $N$ is the number of $(\delta,s)$-supported points. Also, no more than $c\delta^{-2}$ points of $C$ can be cities in a single square $S$. It follows from Claim 5.5 that $N\le A|C|\delta^{-2}\ln(\delta^{-1})/s$, concluding the proof of Lemma 5.3.

\paragraph{Recurrence of bounded degree planar graph limits (Section 5.3).}
Theorem 5.2 follows immediately from the following theorem which gives a quantitative estimate on the growth of the resistance in local limits of bounded degree planar maps.
\begin{theorem}[5.8]
Let $(U,\rho)$ be a local limit of (possibly random) finite planar maps with maximum degree at most $D$. Then, almost surely, there exist a constant $c>0$ and a sequence $\{B_k\}_{k\ge1}$ of subsets of $U$ such that for each $k$ we have
\[
|B_k|\le c^{-1}k,\qquad \Reff(\rho\leftrightarrow U\setminus B_k)\ge c\log k.
\]
In particular, $(U,\rho)$ is almost surely recurrent.
\end{theorem}
We write $B_{\rm euc}(p,r)$ for the Euclidean ball of radius $r$ around a point $p\in\R^2$. As before, for a subset $O\subset\R^2$ and a given circle packing we write $V_O$ for the set of vertices in which the centers of the corresponding circles are in $O$.
\begin{corollary}[5.9]
Let $G$ be a finite simple planar triangulation, and $P$ a circle packing of $G$. Let $\rho$ be a uniform random vertex and $P'$ a dilation and translation of $P$ such that the circle of $\rho$ is a unit circle centered at the origin $\mathbf0$. Then, there exists a universal constant $A>0$ such that in the packing $P'$, for every real $r\ge2$ and integer $s\ge2$
\[
\Prob\left(\forall p\in\R^2\quad\left|V_{B_{\rm euc}(\mathbf0,r)\setminus B_{\rm euc}(p,1/r)}\right|\ge s\right)\le\frac{Ar^2\log r}{s}.
\]
\end{corollary}
\begin{proof}
Apply the Magic Lemma with $\delta=1/r$ and $s=s$, with the centers of circles of $P'$ as the point set $C$. Note that there exists a constant $C>0$ such that for all $w\in V$ the isolation radius of $w$, $\rho_w$, satisfies $\operatorname{rad}(C_w)\le\rho_w\le C\operatorname{rad}(C_w)$ (without appealing to the Ring Lemma).
\end{proof}
\begin{lemma}[5.10]
Let $G$ be a finite simple planar map with maximum degree at most $D$ and let $\rho$ be a uniform random vertex of $G$. Then, there exists a constant $C=C(D)<\infty$ such that for all $k\ge1$,
\begin{multline*}
\Prob\left(\exists B\subseteq V,\ |B|\le Ck,\ \Reff(\rho\leftrightarrow V\setminus B)\ge C^{-1}\log k\right)\\
\ge1-Ck^{-1/3}\log k,
\end{multline*}
where we interpret $\Reff(\rho\leftrightarrow V\setminus B)=\infty$ when $B=V$.
\end{lemma}
\begin{proof}
We first assume that $G$ is a triangulation and consider a circle packing of it where the circle of $\rho$ is a unit circle centered at the origin $\mathbf0$. Applying Corollary 5.9 with $r=k^{1/3},s=k$, we have that with probability at least $1-Ak^{-1/3}\log(k)/3$, there exists $p\in\R^2$ with
\[
\left|V_{B_{\rm euc}(\mathbf0,r)\setminus B_{\rm euc}(p,1/r)}\right|<k.
\]
Now, if $|V_{B_{\rm euc}(p,1/r)}|\le1$, we set $B=V_{B_{\rm euc}(\mathbf0,r)}$. We then have $|B|\le k$ and by applying $\Omega(\log k)$ times Lemma 4.9 together with the series law (Claim 2.24) we get that $\Reff(\rho\leftrightarrow V\setminus B)\ge c\log k$ for some $c=c(D)>0$.
Else, if $|V_{B_{\rm euc}(p,1/r)}|\ge2$ then we take $B=V_{B_{\rm euc}(\mathbf0,r)}\setminus V_{B_{\rm euc}(p,1/r)}$. By the Ring Lemma, there exists a $c'=c'(D)>0$ such that $|p|\ge1+c'$. Since $|V_{B_{\rm euc}(p,1/r)}|\ge2$, we have a vertex in that set with radius at most $r^{-1}$. Therefore, $B_{\rm euc}(p,2/r)$ contains at least one full circle $C_v$. Hence, by scaling and translating such that $C_v=\mathbb U$, we get (again, by Lemma 4.9) that
\[
\Reff\left(V_{B_{\rm euc}(p,2/r)}\leftrightarrow V\setminus V_{B_{\rm euc}(p,c'/2)}\right)\ge c_2\log k,
\]
for some other constant $c_2=c_2(D)>0$. Since $\rho\in V\setminus V_{B_{\rm euc}(p,c'/2)}$ we obtain
\[
\Reff\left(\rho\leftrightarrow V_{B_{\rm euc}(p,2/r)}\right)\ge c_2\log k.
\]
By Lemma 4.9 we also have that $\Reff(\rho\leftrightarrow V\setminus V_{B_{\rm euc}(\mathbf0,r)})\ge c_3\log k$, for some $c_3=c_3(D)>0$. By Claim 2.22 this means that
\begin{align*}
\Prob_\rho(\tau_{V\setminus V_{B_{\rm euc}(\mathbf0,r)}}<\tau_\rho^+)&\le\frac1{c_2\log k},\\
\Prob_\rho(\tau_{V_{B_{\rm euc}(p,2/r)}}<\tau_\rho^+)&\le\frac1{c_3\log k}.
\end{align*}
By the union bound $\Prob_\rho(\tau_{V\setminus B}<\tau_\rho^+)\le2/(\min(c_2,c_3)\log k)$, hence by Claim 2.22 again
$\Reff(\rho\leftrightarrow V\setminus B)\ge\min(c_2,c_3)D^{-1}\log(k)/2$,
concluding the proof when $G$ is a triangulation.

If $G$ is not a triangulation, we would like to add edges to make it a triangulation while making sure that the maximal degree does not increase too much. We also have to ensure that the graph remains simple which may require us to add some additional vertices as well. Let $f$ be a face of $G$ with vertices $v_1,\ldots,v_k$. Suppose first that there are no edges between non-consecutive vertices of the face. In this case, we draw the edges in a zig-zag fashion, as in Fig. 5.3.
In the case where there are edges between non-consecutive vertices of the face exist, we draw a cycle $u_1,\ldots,u_k$ inside $f$. Then, we connect $u_i$ to $v_i$ and $v_{i+1}$ for each $i<k$ and $u_k$ to $v_k$ and $v_1$. Finally, we triangulate the inner face created by the new cycle by zig-zagging as in the previous case (see Fig. 5.4).
Since each vertex of the original graph is a member of at most $D$ faces and for each face at most $2$ edges are added, the maximal degree of the resulting graph is at most $3D$. Similarly, the number of vertices in the resulting graph is at most $D$ times the number of vertices in the original graph hence the probability of a random vertex being a vertex of the original graph is at least $D^{-1}$. If this occurs then it is straightforward to see that the existence of a subset of vertices $B$ in the new graph which satisfies the required conditions implies the existence of such a set in the old graph, concluding our reduction to the triangulation case and finishing our proof.
\end{proof}

\begin{proof}[Proof of Theorem 5.8]
Assume that $G_n$ are finite planar maps with maximum degree at most $D$ such that $G_n\xrightarrow{\mathrm{loc}}(U,\rho)$. If $\{G_n\}$ are not simple graphs we erase self-loops and merge parallel edges into a single edge to obtain the sequence $\{G_n'\}$. It is immediate that $G_n'\xrightarrow{\mathrm{loc}}(U',\rho')$ where $(U',\rho')$ is distributed as $(U,\rho)$ after removing from $U$ all loops and merging parallel edges into a single edge. Since the maximum degree is bounded, $U'$ is recurrent if and only if $U$ is recurrent. Thus we may assume that $G_n$ are simple graphs so the previous estimates may be used.
Denote by $\mathcal A_k$ the event
\[
\mathcal A_k=\{\exists B\subseteq U,\ |B|\le Ck,\ \Reff(\rho\leftrightarrow V\setminus B)\ge c\log k\},
\]
where $C=C(D)<\infty$ is the constant from Lemma 5.10. Therefore $\Prob(\mathcal A_k^c)\le c^{-1}k^{-1/3}\log k$. Looking at the sequence $\{\mathcal A_{2^j}\}_{j\ge1}$, by Borel--Cantelli, almost surely there exists $j_0$ such that for all $j\ge j_0$ the event $\mathcal A_{2^j}$ holds. Thus we have proved the required assertion for $k$ which is a power of $2$.
To prove this for all $k$ sufficiently large, let $B_{2^j}$ be the set guaranteed to exist in the definition of $\mathcal A_{2^j}$, and take $B_k=B_{2^j}$ for the unique $j$ for which $2^j\le k<2^{j+1}$. It is immediate that these sets satisfy the assertion of the theorem, concluding our proof.
\end{proof}

\paragraph{Referenced annular estimate: Lemma 4.9 and its proof (Chapter 4).}
Denote the circle packing $P=\{C_v\}_{v\in V}$ where $V$ is the vertex set of $G$ and $C_v$ denotes the circle corresponding to the vertex $v$. Write $\Delta$ for the maximum degree of $G$ and fix a vertex $v_0$. By scaling and translating we may assume that $C_{v_0}$ is a radius $1$ circle around the origin. For a real number $R>0$, let $V_R=V_{B(0,R)}$ denote set of vertices $v$ for which $\operatorname{cent}(v)$ is in the Euclidean ball of radius $R$ around the origin.
\begin{lemma}[4.9]
There exist $C=C(\Delta)>1$ and $c=c(\Delta)>0$ such that for every $R\ge1$ we have
\begin{enumerate}[label=(\roman*)]
\item There are no edges between $V_R$ and $V\setminus V_{CR}$, and
\item $\Reff(V_R\leftrightarrow V\setminus V_{CR})\ge c$.
\end{enumerate}
\end{lemma}
\begin{proof}
We begin with part (i). For every $v\in V_R$ it holds that $\operatorname{rad}(v)\le R$ since $C_{v_0}$ is centered at the origin. By the Ring Lemma (Lemma 4.2), there exists $A=A(\Delta)$ such that $\operatorname{rad}(u)\le AR$ for every $u\sim v$, and therefore $|\operatorname{cent}(u)|\le(A+2)R$. Hence (i) holds with $C=A+2$.
To prove part (ii) we define
\[
h(v)=\begin{cases}
0&v\in V_R,\\1&v\in V\setminus V_{CR},\\
\dfrac{|\operatorname{cent}(v)|-R}{(C-1)R}&\text{otherwise}.
\end{cases}
\]
Recall from Lemma 2.32 that $\Reff(V_R\leftrightarrow V\setminus V_{CR})\ge\mathcal E(h)^{-1}$. By the triangle inequality, for an edge $\{x,y\}$ with both endpoints in $V_{CR}\setminus V_R$ we have
\[
|h(x)-h(y)|\le\frac{|\operatorname{cent}(x)-\operatorname{cent}(y)|}{(C-1)R}
=\frac{\operatorname{rad}(x)+\operatorname{rad}(y)}{(C-1)R},
\]
and it is straightforward to check that the same bound holds also when one of the edge's endpoints is in $V_R$ or $V\setminus V_{CR}$. Thus, using the Ring Lemma's (Lemma 4.2) constant $A=A(\Delta)$ from part (i),
\begin{align*}
\mathcal E(h)
&\le\sum_{x\in V_{CR}\setminus V_R}\sum_{y:y\sim x}
\frac{((A+1)\operatorname{rad}(x))^2}{(C-1)^2R^2}\\
&\le\frac{\Delta(A+1)^2}{\pi(C-1)^2R^2}\sum_{x\in V_{CR}\setminus V_R}\operatorname{area}(C_x),
\end{align*}
where $\operatorname{area}(C_x)$ is the area that $C_x$ encloses (that is, $\pi\operatorname{rad}(x)^2$). We have that $\sum_x\operatorname{area}(C_x)\le\operatorname{area}(B(\mathbf0,2CR))=4\pi C^2R^2$, hence if $C=A+2$, then $\mathcal E(h)\le4\Delta C^2$, and the result follows for $c=(4\Delta C^2)^{-1}$.
\end{proof}
\subsection*{整理说明（非作者正文）}
题面保留作者一般平面图的版本，证明包含随机多尺度方格的 Magic Lemma、圆堆积转译、电阻估计、非三角剖分的有界度增补、局部极限与 Borel--Cantelli。\S5.1 的动机和例子不收入；必要的局部收敛定义保留，末句图球记号为据原定义整理的说明。圆堆积存在定理、电阻和首次命中概率关系、Rayleigh 单调性及 Borel--Cantelli 为先修；本题并非直接引用常返性定理。

难点在于把均匀随机根转换为几何点集的支持点计数，构造小顶点集外的对数电阻障碍，再传递到局部极限。原图5.3、5.4的锯齿添边和内圈构造保留原图号及下方原页链接，文字构造和度数控制论证完整；图像未取得本地副本。原文 $c_2,c_3$ 在命中概率式中的顺序、最后事件中的 $U/V$ 记号、部分冠词均沿原文，不作数学重审。网页和 PDF 文本转录，非原件下载。
\par 原图与正文：\url{https://link.springer.com/chapter/10.1007/978-3-030-27968-4_5#Fig3}\quad\url{https://link.springer.com/chapter/10.1007/978-3-030-27968-4_5#Fig4}
\subsection*{可修改的简短标注（非作者正文）}
$c$：随机根平面图局部极限的常返性。

$a$：局部弱收敛；圆堆积；随机尺度方格；质量流计数；均匀随机根；有效电阻；命中概率；Borel--Cantelli。

\texttt{cross\_theory\_status = yes}。几何支持点计数经均匀随机根转成高概率的小集电阻障碍；电阻与游走命中概率的对应及局部收敛把有限图估计变为无限极限的常返性。位置：Lemmas 5.3、5.10 与 Theorem 5.8。
标签为非穷尽、可修改初稿。
\subsection*{许可与修改记录}
原作者及作品见来源栏；来源采用 CC BY 4.0。本题证摘录和附加整理沿用该许可。2026-09-28：增加题号、来源定位、独立排版、整理说明和初稿标注；数学内容的取材范围及排版变化见整理说明。
\end{document}
